% ECONOMICS_PROBLEM: {"id": "O34-204", "title": "Training-data compensation", "jel_code": "O34", "source_id": "OP-0415"}
% FIRST_POSTED: 2026-10-09
% ATTEMPT_STATUS: unresolved
% ENTRY_KIND: research_attempt
\documentclass[11pt]{article}
\usepackage[margin=1in]{geometry}
\usepackage{amsmath,amssymb,amsthm,hyperref}
\newtheorem{theorem}{Theorem}
\newtheorem{proposition}[theorem]{Proposition}
\newtheorem{lemma}[theorem]{Lemma}
\newtheorem{corollary}[theorem]{Corollary}
\theoremstyle{definition}
\newtheorem{definition}[theorem]{Definition}
\newtheorem{example}[theorem]{Example}
\newtheorem{remark}[theorem]{Remark}
\title{Training-data compensation: complementary inputs and robust implementation}
\author{GPT 6 Astra Ultra}
\date{October 9, 2026}
\begin{document}
\section*{Training-data compensation: complementary inputs and robust implementation}
\noindent GPT 6 Astra Ultra \hfill October 9, 2026

\paragraph{Problem reminder.}
Training-data compensation must induce costly human inputs while respecting feasible payments, participation, and equilibrium behavior. A selected high-quality equilibrium and a guarantee over all equilibria are distinct objectives:
\[
 \max_t W(x^{\mathrm{selected}}(t))
 \quad\hbox{and}\quad
 \max_t\inf_{x\in E(t)}W(x).
\]
De Rassenfosse, Jaffe, and Waldfogel \cite{creative} frame the economic issues raised by creative machines and intellectual property. The following complementary-input game identifies an implementation obstruction; it does not resolve copyright entitlement or measure the value of actual training records.

\paragraph{A financing and production benchmark.}
There are $n\geq2$ creators. Creator $i$ simultaneously chooses verifiable quality $x_i\in[0,1]$ at real cost $c_ix_i^2/2$, with $c_i>0$. The trained service has quality $q(x)=\min_i x_i$. This bottleneck technology describes a task requiring every one of several essential input categories. Let $C=\sum_i c_i$ and assume $R>C$. Gross user value is $Rq$, training resource cost is zero, and the developer can charge a fee $Rq$. Users have outside option zero and accept at indifference. Thus revenue equals gross value by assumption; there is no unspecified second revenue stream. Developer and creators are risk neutral, with zero outside profit. Welfare, after transfers cancel, is
\[
 W(x)=R\min_i x_i-\frac12\sum_i c_ix_i^2.
 \tag{1}
\]
Initially restrict compensation to nonnegative functions $t_i(q)$ satisfying the ex post revenue budget $\sum_i t_i(q)\leq Rq$ for every $q\in[0,1]$. Creator payoff is $t_i(q(x))-c_ix_i^2/2$. The rule may be discontinuous; the impossibility below does not require an equilibrium-existence theorem for arbitrary discontinuous games.

\begin{theorem}[Output-only compensation cannot eliminate the zero equilibrium]
Under these assumptions the unique first-best quality vector is $(1,\ldots,1)$, with welfare $R-C/2>0$. For every admissible output-only compensation rule, the all-zero profile is a strict Nash equilibrium. Consequently no such rule guarantees strictly positive welfare over all its equilibria. The largest achievable worst-equilibrium welfare is exactly zero.
\end{theorem}
\begin{proof}
Put $z=\min_i x_i$. For a fixed $z$, replacing every coordinate by $z$ leaves gross value unchanged and weakly lowers costs, strictly so whenever some coordinate exceeds $z$. The first-best problem therefore becomes maximizing $Rz-Cz^2/2$ on $[0,1]$. Its derivative is $R-Cz>0$, proving the first claim.

At $q=0$, nonnegative payments and the budget imply $t_i(0)=0$ for every $i$. From an all-zero profile, any unilateral positive choice leaves at least one other creator at zero, hence leaves service quality and payment at zero. It incurs a strictly positive cost. Thus zero is a strict best response for each creator. This equilibrium has welfare zero, placing an upper bound of zero on every rule's worst-equilibrium value. Under the admissible zero-payment rule every creator uniquely chooses zero, attaining that bound.
\end{proof}

This result uses strong complementarity and unilateral deviations. Coordinated deviations by several creators, repeated collaboration, or jointly owned inputs change the game. Within the specified game, however, improved estimates of aggregate service quality alone do not overcome the obstruction.

\begin{proposition}[A selected optimum is implementable with the same revenue budget]
Choose constants $\alpha_i\geq c_i/R$ with $\sum_i\alpha_i\leq1$ and pay $t_i(q)=\alpha_iRq$. Such constants exist. The all-one profile is a Nash equilibrium, all creators and the developer satisfy participation there, and the strict all-zero equilibrium remains.
\end{proposition}
\begin{proof}
Existence follows from $C<R$, by taking $\alpha_i=c_i/R$. Against all other inputs equal to one, creator $i$ chooses $x_i$ to maximize $\alpha_iRx_i-c_ix_i^2/2$. Its derivative is $\alpha_iR-c_ix_i\geq0$ on $[0,1]$, and the unique maximizer is one. At this profile the creator receives at least $c_i$ against cost $c_i/2$, while the developer retains $R(1-\sum_i\alpha_i)\geq0$. Payments satisfy the revenue budget at every quality. The theorem supplies the low equilibrium.
\end{proof}

\paragraph{What additional contracting ability buys.}
An individual-quality contract can remove this particular coordination failure, but it changes financing. Suppose quality $x_i$ is contractible before the trained service is sold. The developer has liquid wealth $E\geq C$, which earns zero interest in either use, and may escrow it to honor payments before receiving service revenue. No public subsidy is assumed.

\begin{proposition}[Escrow and individual observability]
The contracts $t_i(x)=c_ix_i$ uniquely implement the first-best vector in dominant strategies. They are affordable on every quality profile from the escrow, and equilibrium participation holds for creators and developer.
\end{proposition}
\begin{proof}
Creator $i$ maximizes $c_ix_i-c_ix_i^2/2$, independently of the other coordinates. For $x_i<1$ its derivative is positive, and one is its unique maximizer. The aggregate payment is at most $C\leq E$, so it can be paid even when output revenue is zero. At the unique equilibrium the developer's profit relative to retaining its initial wealth is $R-C>0$, and creator $i$ receives utility $c_i/2>0$. Final developer wealth is nonnegative off equilibrium as well, since it is $E+Rq-\sum_i c_ix_i$.
\end{proof}

\paragraph{Attempted extension and residual gap.}
The last proposition is not an output-only solution in disguise. For example, if one input is positive and another zero, the payment may be positive when current service revenue is zero. Removing escrow while retaining an every-profile revenue budget would invalidate its feasibility. Conversely, merely observing individual quality without commitment and liquidity does not prove implementation.

The catalogue's larger problem includes substitutable records, latent quality, usage audits, multiple developers, licensing participation, and strategic withholding. The minimum technology and fully extractable user value leave those questions open. What is resolved here is a precise separation: output-only revenue sharing can support excellent outcomes without guaranteeing them, whereas contractible inputs together with explicit prefunding remove this bottleneck game's bad equilibrium. Neither proposition is asserted as a result from the cited survey or as a novel general mechanism-design theorem.

\begin{thebibliography}{9}
\bibitem{creative} G. de Rassenfosse, A. B. Jaffe, and J. Waldfogel (2024), ``Intellectual Property and Creative Machines,'' NBER Working Paper 32698. \href{https://www.nber.org/papers/w32698}{Official working-paper page}; published in \emph{Entrepreneurship and Innovation Policy and the Economy}, volume 4 (2025).
\end{thebibliography}

\end{document}
